\documentclass[12pt,a4paper]{ctexart}

% ===== 基础宏包 =====
\usepackage[top=2.5cm, bottom=2.5cm, left=2.8cm, right=2.8cm]{geometry}
\usepackage{graphicx}
\usepackage{booktabs}
\usepackage{longtable}
\usepackage{array}
\usepackage{tabularx}
\usepackage{enumitem}
\usepackage{xcolor}
\usepackage{hyperref}
\usepackage{xurl}
\usepackage{fancyhdr}
\usepackage{titlesec}
\usepackage{caption}
\usepackage{multirow}
\usepackage{makecell}
\usepackage{amssymb}
\usepackage{float}        % 提供 [H] 强制表格定位，防止浮动跳位

% ===== 颜色定义 =====
\definecolor{primary}{RGB}{41,65,122}
\definecolor{accent}{RGB}{200,60,60}
\definecolor{lightgray}{RGB}{245,245,245}
\definecolor{boxbg}{RGB}{240,244,252}

% ===== 超链接样式 =====
\hypersetup{
    colorlinks=true,
    linkcolor=primary,
    urlcolor=primary,
    citecolor=primary
}

% ===== 标题样式 =====
\titleformat{\section}{\Large\bfseries\color{primary}}{\thesection}{1em}{}[\vspace{2pt}\hrule]
\titleformat{\subsection}{\large\bfseries}{\thesubsection}{1em}{}
\titleformat{\subsubsection}{\normalsize\bfseries}{\thesubsubsection}{1em}{}

% ===== 页眉页脚 =====
\pagestyle{fancy}
\fancyhf{}
\fancyhead[L]{\small\color{gray}BitDance 舞蹈学习与约练平台}
\fancyhead[R]{\small\color{gray}产品立项建议书}
\fancyfoot[C]{\thepage}
\renewcommand{\headrulewidth}{0.4pt}
\setlength{\headheight}{14pt}

% ===== 列表样式 =====
\setlist[itemize]{leftmargin=2em, itemsep=2pt, parsep=0pt}
\setlist[enumerate]{leftmargin=2em, itemsep=2pt, parsep=0pt}

% ===== 表格列类型 =====
\newcolumntype{L}[1]{>{\raggedright\arraybackslash}p{#1}}
\newcolumntype{C}[1]{>{\centering\arraybackslash}p{#1}}

% ===== 场景框 =====
\usepackage{tcolorbox}
\tcbuselibrary{skins,breakable}
\newtcolorbox{scenariobox}[1]{
    colback=boxbg, colframe=primary, boxrule=0.5pt,
    fonttitle=\bfseries, title={#1},
    left=8pt, right=8pt, top=4pt, bottom=4pt,
    breakable
}

\begin{document}

% ===== 封面 =====
\begin{titlepage}
\centering
\vspace*{4cm}
{\Huge\bfseries\color{primary} BitDance\par}
\vspace{0.8cm}
{\LARGE 舞蹈学习与约练平台\par}
\vspace{0.5cm}
{\Large 产品立项建议书\par}
\vspace{3cm}
{\large 第五组\par}
\vspace{2cm}
{\large \today\par}
\vfill
\end{titlepage}

\newpage
\tableofcontents
\newpage

% ============================================================
\section{项目概述}
% ============================================================

\subsection{项目名称与定位}

BitDance 是一款面向舞蹈爱好者的学习与约练平台，围绕\textbf{找舞室、选课程、约搭子、持续练习、成长记录}五个核心环节，提供比现有地图平台、点评平台和短视频平台更完整的决策与学习支持。

平台的核心目标是围绕舞蹈学习者的真实路径建立一条清晰的服务链路：
\begin{enumerate}
    \item 帮助用户找到适合自己的舞室和课程；
    \item 帮助用户找到同伴并建立练习节奏；
    \item 通过成长记录和内容沉淀提高长期留存。
\end{enumerate}

\subsection{立项结论}

舞蹈学习兼具兴趣、健身、社交和自我表达属性，是长期行为而非一次性消费。当前用户在找舞室、比课程、看评价、找搭子和坚持练习的过程中需在多个平台间来回切换，信息分散、判断困难、试错成本高。本项目的价值在于将这些分散需求整合为一个围绕舞蹈学习路径设计的垂直平台。该方向技术边界清晰、实现路径明确、商业模式具备现实空间。

% ============================================================
\section{市场分析}
% ============================================================

\subsection{宏观市场环境}

舞蹈学习同时叠加兴趣消费、健身需求、社交需求和自我表达需求，并非短期热点，而是年轻用户中持续存在的服务品类。以下从消费趋势、产业规模、政策支撑和内容生态四个维度论证市场基础。

\subsubsection{消费趋势：文化娱乐支出持续增长}

国家统计局发布的《2025 年居民收入和消费支出情况》显示，2025 年全国居民人均教育文化娱乐消费支出为 3489 元，同比增长 9.4\%\textsuperscript{[1]}。舞蹈作为兼具兴趣、健身和线下体验属性的细分品类，直接受益于这一增长趋势。与纯线上娱乐消费不同，舞蹈学习涉及线下课程付费、器材购置和活动参与等多重消费环节，客单价和复购率均高于一般文娱消费。

\subsubsection{产业规模：街舞培训已形成成熟网络}

中国舞蹈家协会街舞委员会发布的《中国街舞产业发展研究报告》\textsuperscript{[3]}显示，截至 2023 年：

\begin{table}[H]
\centering
\caption{中国街舞产业核心数据}
\begin{tabularx}{\textwidth}{L{4cm}X}
\toprule
\textbf{指标} & \textbf{数据} \\
\midrule
CHUC 体系覆盖城市 & 596 个 \\
街舞教育培训机构 & 近 10,000 家 \\
从业者规模 & 近 300 万人 \\
年学习人次 & 超 1,000 万 \\
核心目标群体 & Z 世代（1995--2009 年出生） \\
\bottomrule
\end{tabularx}
\end{table}

这意味着街舞、Hiphop、Urban、Locking、Breaking 等面向青年群体的舞种已具备稳定的供给侧基础设施（舞室、师资、课程体系），平台对接的不是零散个人兴趣，而是已经成型的产业网络。

\subsubsection{政策支撑：街舞从圈层走向主流}

国家体育总局办公厅在关于高校共建奥运会新兴项目后备人才培养体系的答复中明确提到，霹雳舞（Breaking）已列入 2024 年巴黎奥运会正式比赛项目，并提出推动新兴项目在高校开展、带动后备人才培养\textsuperscript{[4]}。这标志着街舞不再只是娱乐内容，而在年轻人语境中同时具备\textbf{潮流表达、体育训练和校园参与}三重属性。高校舞社的普及和赛事体系的建立，进一步为平台提供了天然的种子用户池和推广渠道。

\subsubsection{内容生态：短视频平台提供传播基础设施}

第 56 次《中国互联网络发展状况统计报告》\textsuperscript{[6]}显示，截至 2025 年 6 月，我国短视频用户规模已达 10.68 亿人，占网民整体的 95.1\%。第一财经 2024 年文章援引研究指出，全球韩流粉丝规模已达 2.25 亿，核心受众为青年群体\textsuperscript{[5]}。

短视频平台上韩舞翻跳、课堂记录、Workshop 宣发和舞室种草的高频传播形态表明：用户不只是被动观看，而是围绕作品模仿、跟练打卡、同伴结识和线下体验持续产生后续需求。\textbf{现有内容平台擅长激发兴趣，但不负责将内容热度转化为课程决策、搭子关系和长期训练记录}——这正是 BitDance 的切入点。

\subsection{目标市场定位}

本项目不追求覆盖所有兴趣培训市场，而是聚焦\textbf{年轻化城市舞蹈}这一细分场景。与书法、乐器等偏静态学习相比，街舞、Hiphop、韩舞、Urban 等舞种更依赖线下空间、同伴关系、视频传播和阶段展示，更容易形成"线上发现 → 线下消费 → 线上复盘 → 反哺线下"的循环。

收入来源不只是常规课程付费，还包括更高频和更有传播性的周边场景：

\begin{itemize}
    \item \textbf{课程层}：长期班、体验课、私教课
    \item \textbf{活动层}：Workshop、公开课、假日集训、成品舞拍摄
    \item \textbf{社交层}：约练拼课、舞室推广
    \item \textbf{内容层}：课堂记录传播带来的自然获客
\end{itemize}

\subsection{市场空白分析}

当前市场上不缺单点工具，缺的是一套围绕舞蹈学习路径持续服务的系统。具体体现为四个维度的空白：

\begin{table}[H]
\centering
\caption{市场空白与现有方案对比}
\begin{tabularx}{\textwidth}{L{3cm}L{4cm}X}
\toprule
\textbf{空白维度} & \textbf{用户真实需求} & \textbf{现有方案的不足} \\
\midrule
垂直搜索 & 按舞种、课程适配度、老师风格精准搜索 & 美团/地图仅支持泛品类搜索，无舞蹈专属筛选维度 \\
\midrule
结构化评价 & 了解"零基础能否跟上""老师是否耐心""课程节奏如何" & 大众点评以环境和服务为主，评价维度与舞蹈学习决策错位 \\
\midrule
约练与社交 & 找到水平匹配的练舞搭子，约定时间地点一起练 & 微信群/QQ 群效率低、无水平匹配、响应率不可控 \\
\midrule
成长记录 & 记录训练时长、课程进度、阶段作品，看到自己的进步 & 无现有平台提供面向舞蹈学习者的长期成长档案 \\
\bottomrule
\end{tabularx}
\end{table}

\subsection{竞品分析}

\subsubsection{竞品分类与对比}

现有竞品按功能定位分为三类：

\begin{table}[H]
\centering
\caption{竞品详细对比}
\begin{tabularx}{\textwidth}{L{2cm}L{2.8cm}L{3cm}X}
\toprule
\textbf{竞品类型} & \textbf{代表产品} & \textbf{优势} & \textbf{不足} \\
\midrule
舞室自有系统 & 各舞室小程序、公众号、私域群 & 信息更新快；课表管理直接；成交路径短；老带新效率高 & 信息孤岛，缺少横向比较；不适合帮新用户做跨舞室决策；无法承载成长记录 \\
\midrule
本地生活平台 & 大众点评、美团、高德/百度地图 & 商家覆盖广；位置/价格/团购信息完善；用户基数大 & 评价维度偏泛消费（环境、服务），无法回答课程适配和教学质量问题；报名后无后续服务 \\
\midrule
内容与社交平台 & 抖音、小红书、微信群、朋友圈 & 种草内容丰富；作品展示直观；同好发现自然 & 信息碎片化；评价不成体系；约练效率低下；无法形成持续的学习关系 \\
\bottomrule
\end{tabularx}
\end{table}

\subsubsection{BitDance 的差异化定位}

BitDance 的竞争力不在于替代任何一类现有平台，而在于将三类平台各自承担的一小段功能\textbf{重新组织为一条从决策到成长的完整路径}。

\begin{table}[H]
\centering
\caption{BitDance 与现有方案的链路对比}
\begin{tabularx}{\textwidth}{L{2.5cm}X X}
\toprule
\textbf{用户路径环节} & \textbf{现有方案} & \textbf{BitDance} \\
\midrule
找舞室 & 美团/地图搜索 → 泛品类结果 & 按舞种、适合人群、课程强度精准筛选 \\
选课程 & 短视频/小红书种草 → 碎片信息 & 结构化课程详情 + 多舞室横向对比 \\
看评价 & 大众点评 → 环境/服务维度 & 专为舞蹈设计的三维评价体系（舞室/老师/课程） \\
找同伴 & 微信群/私聊 → 低效撮合 & 约练广场 + 智能匹配 + 搭子关系沉淀 \\
坚持训练 & 无平台支持 & 训练打卡 + 成长时间线 + 阶段作品记录 \\
\bottomrule
\end{tabularx}
\end{table}

大平台不会专门为舞蹈学习这一细分领域深挖体验细节。\textbf{谁能把决策、陪伴和成长三段链路打通，谁就更有机会在这个垂直赛道建立壁垒。}

\subsubsection{目标人群意愿验证}

在小红书发布意见征集帖后获得积极反响，用户表现出较强期待。甚至有香港科技大学毕业生在帖下留言表示希望加入共同创业，侧面验证了需求的真实性和市场的想象空间。

% ============================================================
\section{用户与需求分析}
% ============================================================

\subsection{用户画像}

\subsubsection{画像 A：舞蹈入门新手}

\begin{itemize}
    \item \textbf{年龄范围}：18--28 岁，以大学生和初入职场的年轻人为主
    \item \textbf{典型特征}：对舞蹈感兴趣但尚未开始或刚开始学习，信息获取依赖短视频种草和朋友推荐
    \item \textbf{核心痛点}：
    \begin{itemize}
        \item 不知该选哪家舞室、哪个舞种、哪位老师
        \item 担心零基础跟不上课程节奏
        \item 缺少同伴，一个人报名容易放弃
        \item 在多个平台间反复搜索比较，决策成本高
    \end{itemize}
    \item \textbf{核心诉求}：快速找到适合零基础的舞室和课程，获得真实可信的评价参考，最好能找到一起学的搭子
\end{itemize}

\subsubsection{画像 B：持续学舞的进阶用户}

\begin{itemize}
    \item \textbf{年龄范围}：18--35 岁，有半年以上学舞经验
    \item \textbf{典型特征}：已有固定学习的舞种，关注特定老师和 Workshop，有练舞习惯
    \item \textbf{核心痛点}：
    \begin{itemize}
        \item 想找水平相当的练舞搭子，但线下组织效率低
        \item Workshop 信息分散，容易错过报名
        \item 缺少系统化的成长记录和进步可视化
        \item 想尝试新舞种或新舞室，但缺乏针对性参考
    \end{itemize}
    \item \textbf{核心诉求}：高效约练、Workshop 信息聚合、成长记录沉淀、跨舞室课程发现
\end{itemize}

\subsubsection{画像 C：舞室经营者/独立教练}

\begin{itemize}
    \item \textbf{典型特征}：经营舞蹈工作室或以个人身份授课
    \item \textbf{核心痛点}：
    \begin{itemize}
        \item 获客依赖美团和私域，缺少精准触达舞蹈学习者的渠道
        \item 课程信息展示受限于泛平台模板，无法突出教学特色
        \item 缺少对学员画像的了解，难以优化课程设置
    \end{itemize}
    \item \textbf{核心诉求}：精准获客、专业化课程展示、学员管理、品牌建设
\end{itemize}

\subsection{需求层次分析}

基于三类用户画像和业务路径，将产品需求按 KANO 模型与优先级划分如下：

\begin{table}[H]
\centering
\caption{需求优先级矩阵}
\begin{tabularx}{\textwidth}{C{0.8cm}L{1.2cm}L{3cm}X L{2cm}}
\toprule
\textbf{级别} & \textbf{属性} & \textbf{需求} & \textbf{说明} & \textbf{对应模块} \\
\midrule
P0 & 基本型 & 舞室发现与搜索 & 按位置、舞种、价格等条件快速找到舞室 & 舞室搜索 \\
P0 & 基本型 & 课程信息展示 & 结构化展示课程详情，支持多维度对比 & 课程详情 \\
P0 & 基本型 & 结构化评价 & 围绕舞蹈学习场景的专业评价体系 & 评价系统 \\
P0 & 基本型 & 用户账号体系 & 注册、登录、个人资料管理 & 账号模块 \\
\midrule
P1 & 期望型 & 约搭子/约练 & 基于舞种、位置、时间匹配练舞同伴 & 约练社交 \\
P1 & 期望型 & 成长记录与打卡 & 学舞历程记录、训练打卡、数据统计 & 成长档案 \\
P1 & 期望型 & 收藏与关注 & 收藏舞室/课程/老师 & 个人中心 \\
\midrule
P2 & 期望型 & 内容社区 & 分享试听感受、课堂记录、约练日常 & 社区模块 \\
P2 & 期望型 & Workshop 管理 & 发布、报名、候补、签到 & 活动模块 \\
P2 & 期望型 & 商家入驻与管理 & 舞室认领、课表维护、数据看板 & 商家后台 \\
\midrule
P3 & 兴奋型 & 智能推荐 & 基于偏好推荐舞室、课程、搭子 & 推荐系统 \\
P3 & 兴奋型 & 拼课功能 & 凑齐人数联系舞室开课 & 拼课模块 \\
\bottomrule
\end{tabularx}
\end{table}

\subsection{需求分析总结}

通过用户画像与需求层次的交叉分析，BitDance 的产品需求可归纳为以下关键结论：

\begin{enumerate}
    \item \textbf{决策支持是第一优先级}：用户最迫切的需求是在一个平台上完成舞室发现、课程对比和可信评价的闭环——这是平台的核心价值锚点，直接决定用户是否愿意首次使用。
    \item \textbf{学习陪伴是留存关键}：约搭子、约练和成长记录解决的是"报名之后怎么坚持"的问题，直接决定用户生命周期长度。数据表明舞蹈学习中途放弃率极高，同伴机制是对抗流失的核心手段。
    \item \textbf{内容与活动是增长飞轮}：社区内容和 Workshop 活动既能带来新用户，也能提升老用户活跃度，但不应喧宾夺主，需始终服务于学习主线。
    \item \textbf{商家服务是商业化基础}：舞室和教练端的入驻与管理能力是未来变现的核心支撑，但前期以轻量接入为主，避免过早增加复杂度。
\end{enumerate}

\subsection{核心用户场景}

以下四个场景分别覆盖三类用户画像的核心使用路径：

\begin{scenariobox}{场景一：新手找舞室决策}
\textbf{用户}：小李，大二学生，想学韩舞，零基础。

\textbf{场景描述}：小李在短视频上刷到韩舞翻跳视频后想学舞，但不知道学校附近哪家舞室有韩舞课、价格多少、是否适合零基础。

\textbf{当前痛点}：需要在美团搜位置、抖音看风格、小红书看评价、微信群问口碑，来回切换 4 个以上平台。

\textbf{期望体验}：打开 BitDance，定位到学校附近，筛选"韩舞 + 零基础友好"，看到 3 家舞室的课程对比，阅读结构化评价后收藏心仪课程并预约试听。
\end{scenariobox}

\begin{scenariobox}{场景二：进阶用户约练}
\textbf{用户}：小王，上班族，学 Hiphop 一年，想找人周末一起练舞。

\textbf{场景描述}：小王工作日上课、周末想找人一起复习和练习，但舞室微信群里不好意思频繁发约练信息。

\textbf{当前痛点}：群里发约练信息响应率低，且不知对方水平是否匹配。

\textbf{期望体验}：在 BitDance 发布约练信息（Hiphop/中级/周六下午/某舞室），平台推荐匹配的舞友，双方确认后即可约练。
\end{scenariobox}

\begin{scenariobox}{场景三：Workshop 报名}
\textbf{用户}：小张，Urban 舞者，关注某位抖音舞蹈博主。

\textbf{场景描述}：某知名老师要来本城市开 Workshop，小张想第一时间报名。

\textbf{当前痛点}：Workshop 信息散落在不同舞室的公众号和朋友圈，经常错过报名时间，且不确定场地和时间。

\textbf{期望体验}：在 BitDance 收到 Workshop 上线通知，查看详情（时间、地点、价格、老师介绍、往期评价），一键报名并完成支付。
\end{scenariobox}

\begin{scenariobox}{场景四：舞室经营者获客}
\textbf{用户}：陈老师，经营一家主打 Jazz 和韩舞的小型舞室，开业半年。

\textbf{当前痛点}：美团上同类商家竞价激烈，短视频运营成本高，获客效率低。

\textbf{期望体验}：在 BitDance 认领舞室页面，完善课程信息和教练介绍，通过平台的精准匹配触达真正想学韩舞的新手用户。
\end{scenariobox}

% ============================================================
\section{产品方案与功能设计}
% ============================================================

\subsection{体验设计理念}

\textbf{旧体验}：用户在美团搜舞室 → 短视频看风格 → 点评平台看评价 → 微信群问口碑 → 报名后无后续支持。整个过程缺少统一判断标准和持续服务。

\textbf{新体验}：用户在一个平台完成舞室搜索、课程对比、结构化评价、约搭子、训练打卡和成长记录。核心价值不在某个单独功能，而在于将原本割裂的\textbf{决策、陪伴、成长}串成一条连续路径。

\subsection{功能架构总览}

平台功能围绕"\textbf{搜索决策 → 学习陪伴 → 成长沉淀}"的用户主路径展开：

\begin{table}[H]
\centering
\caption{功能模块总览}
\begin{tabularx}{\textwidth}{C{1cm}L{3cm}X C{1cm}}
\toprule
\textbf{编号} & \textbf{模块} & \textbf{核心职责} & \textbf{级别} \\
\midrule
M1 & 舞室与课程 & 搜索、筛选、详情展示、课程信息管理 & P0 \\
M2 & 评价系统 & 结构化评价、评分体系、评价风控 & P0 \\
M3 & 用户账号 & 注册登录、个人资料、偏好设置 & P0 \\
M4 & 约练社交 & 约搭子、约练发布、匹配推荐 & P1 \\
M5 & 成长档案 & 训练打卡、学习记录、数据可视化 & P1 \\
M6 & 社区与活动 & 内容分享、Workshop 管理、活动发布 & P2 \\
M7 & 商家管理 & 舞室入驻、课表维护、教练管理、数据看板 & P2 \\
\bottomrule
\end{tabularx}
\end{table}

\subsection{M1：舞室与课程模块}

\subsubsection{功能清单}

\begin{table}[H]
\centering
\begin{tabularx}{\textwidth}{C{1.5cm}L{2.5cm}X C{0.8cm}}
\toprule
\textbf{编号} & \textbf{功能} & \textbf{描述} & \textbf{级别} \\
\midrule
M1-F01 & 附近舞室搜索 & 基于定位展示附近舞室，支持地图/列表切换 & P0 \\
M1-F02 & 多维度筛选 & 按舞种、价格、距离、时段、适合人群筛选 & P0 \\
M1-F03 & 舞室详情页 & 基础信息、环境照片、主打舞种、课表、导航 & P0 \\
M1-F04 & 课程详情页 & 名称、舞种、难度、价格、老师、训练强度 & P0 \\
M1-F05 & 老师详情页 & 擅长舞种、教学风格、评价、可预约课程 & P0 \\
M1-F06 & 收藏功能 & 收藏舞室/课程/老师 & P1 \\
M1-F07 & 课表查看 & 按日/周视图展示课程安排 & P1 \\
M1-F08 & 试听预约 & 在线预约试听，舞室收到通知 & P1 \\
M1-F09 & 舞室对比 & 2--3 家舞室多维度并排对比 & P2 \\
M1-F10 & 智能推荐 & 根据偏好推荐匹配舞室和课程 & P3 \\
\bottomrule
\end{tabularx}
\end{table}

\subsubsection{核心用例（UC-01）：用户搜索附近舞室}

\begin{itemize}
    \item \textbf{前置条件}：用户已登录或以游客身份进入，已授权地理位置
    \item \textbf{基本流程}：
    \begin{enumerate}
        \item 用户进入首页，系统获取当前位置
        \item 系统展示附近舞室列表（默认按距离排序）
        \item 用户切换地图模式或设置筛选条件（如：韩舞、5km 内、200 元以下）
        \item 系统刷新列表，用户点击舞室卡片进入详情页
    \end{enumerate}
    \item \textbf{异常流程}：未授权定位时提示手动输入地址；无结果时提示扩大范围
\end{itemize}

\subsubsection{业务规则}

\begin{itemize}
    \item 搜索默认范围 5km，最大可扩展至 50km
    \item 舞种采用二级分类：一级为大类（街舞、流行舞、中国舞等），二级为具体舞种（Hiphop、Jazz、韩舞、Urban、Locking、Breaking 等）
    \item 价格筛选支持区间选择：0--50、50--100、100--200、200+ 元/节
    \item 课程难度分四级：零基础友好、初级、中级、高级
    \item 数据来源三层：平台初始录入 → 商家认领完善 → 用户纠错补充
\end{itemize}

\subsection{M2：评价系统模块}

\subsubsection{功能清单}

\begin{table}[H]
\centering
\begin{tabularx}{\textwidth}{C{1.5cm}L{2.5cm}X C{0.8cm}}
\toprule
\textbf{编号} & \textbf{功能} & \textbf{描述} & \textbf{级别} \\
\midrule
M2-F01 & 舞室评价 & 多维度结构化评价（交通、卫生、场地、氛围） & P0 \\
M2-F02 & 老师评价 & 结构化评价（耐心度、纠错质量、讲解清晰度） & P0 \\
M2-F03 & 课程评价 & 结构化评价（上手难度、节奏、强度、收获） & P0 \\
M2-F04 & 评价权重分层 & 根据报名/签到/核销状态区分权重 & P0 \\
M2-F05 & 评分聚合展示 & 综合评分与各维度雷达图 & P1 \\
M2-F06 & 图文/视频评价 & 支持附加照片或短视频 & P1 \\
M2-F07 & 评价风控 & 识别异常评价，降权/折叠/人工复核 & P1 \\
M2-F08 & 商家申诉 & 不实评价申诉，进入人工审核 & P2 \\
\bottomrule
\end{tabularx}
\end{table}

\subsubsection{评价维度设计}

评价体系按三个评价对象分别设计维度，每维度采用 1--5 分制：

\begin{table}[H]
\centering
\caption{三维评价体系}
\begin{tabularx}{\textwidth}{L{1.5cm}L{2.5cm}X}
\toprule
\textbf{对象} & \textbf{维度} & \textbf{说明} \\
\midrule
\multirow{4}{*}{舞室} & 交通便利度 & 到达是否方便（1=很不方便，5=非常方便） \\
 & 环境卫生 & 整洁程度 \\
 & 场地条件 & 镜面、地板、音响等硬件 \\
 & 整体氛围 & 学习和社交氛围 \\
\midrule
\multirow{4}{*}{老师} & 耐心程度 & 对新手的耐心 \\
 & 纠错质量 & 能否准确纠正动作问题 \\
 & 讲解清晰度 & 动作拆解是否清楚 \\
 & 零基础友好 & 是否照顾零基础学员 \\
\midrule
\multirow{4}{*}{课程} & 上手难度 & 对学员能力的要求 \\
 & 节奏合理性 & 教学进度是否适中 \\
 & 练习强度 & 体力消耗程度 \\
 & 实际收获 & 课后能力提升感受 \\
\bottomrule
\end{tabularx}
\end{table}

\subsubsection{评价权重与风控规则}

\begin{itemize}
    \item \textbf{高权重}：通过平台完成报名/签到/核销的用户，标记为"已验证"，权重最高
    \item \textbf{普通权重}：注册用户自主提交，经基础审核后展示
    \item \textbf{低权重/折叠}：新注册账号（<7 天）评价、短时间大量同质化评价、高度相似文案、单店评分异常波动
    \item 舞室不可直接删除评价，仅可通过申诉渠道提交审核
\end{itemize}

\subsection{M3：用户账号模块}

\begin{table}[H]
\centering
\begin{tabularx}{\textwidth}{L{3cm}X C{0.8cm}}
\toprule
\textbf{功能} & \textbf{描述} & \textbf{级别} \\
\midrule
手机号注册/登录 & 手机号 + 验证码 & P0 \\
第三方登录 & 微信授权登录 & P0 \\
个人资料管理 & 昵称、头像、性别、简介 & P0 \\
舞蹈偏好设置 & 感兴趣的舞种、学习水平、目标 & P1 \\
隐私设置 & 控制个人信息可见性 & P1 \\
消息通知 & 系统通知、约练回复、评价互动推送 & P1 \\
社交账号展示 & 个人主页展示其他平台账号 & P2 \\
\bottomrule
\end{tabularx}
\end{table}

\subsubsection{用户角色设计}

\begin{table}[H]
\centering
\begin{tabularx}{\textwidth}{L{2.5cm}X}
\toprule
\textbf{角色} & \textbf{权限} \\
\midrule
游客 & 浏览舞室/课程信息，不可评价、收藏、约练 \\
普通用户 & 搜索、评价、收藏、约练、打卡、社区 \\
舞室管理员 & 管理舞室信息、课表、回复评价、数据看板 \\
教练 & 管理个人主页、查看评价、发布可约课程 \\
平台管理员 & 内容审核、评价风控、数据管理、商家审核 \\
\bottomrule
\end{tabularx}
\end{table}

\subsection{M4：约练社交模块}

\begin{table}[H]
\centering
\begin{tabularx}{\textwidth}{L{3cm}X C{0.8cm}}
\toprule
\textbf{功能} & \textbf{描述} & \textbf{级别} \\
\midrule
发布约练 & 发布舞种、时间、地点、人数、水平要求 & P1 \\
约练广场 & 展示附近/同城约练信息，支持筛选 & P1 \\
约练响应 & 响应他人邀请，双方确认后建立约练关系 & P1 \\
约练匹配推荐 & 根据舞种、水平、位置推荐搭子 & P2 \\
搭子关系 & 互相约练后可互加为搭子 & P2 \\
约练评价 & 约练后双方评价（守时、友好、水平匹配） & P2 \\
拼课发起 & 凑齐人数联系舞室开课 & P3 \\
\bottomrule
\end{tabularx}
\end{table}

\subsubsection{核心用例（UC-03）：发布约练并匹配搭子}

\begin{itemize}
    \item \textbf{前置条件}：用户已登录，已设置舞蹈偏好
    \item \textbf{基本流程}：
    \begin{enumerate}
        \item 用户进入约练广场，点击"发布约练"
        \item 填写约练信息：舞种（Hiphop）、时间（周六 14:00--16:00）、地点（某舞室）、期望人数（2--4 人）、水平要求（中级）
        \item 系统发布至约练广场，同时向符合条件的用户推送通知
        \item 其他用户点击"我要参加"，发起者确认接受
        \item 双方约练关系建立，系统发送时间地点提醒
    \end{enumerate}
    \item \textbf{异常流程}：无人响应则到期自动关闭；发起者取消则通知已响应用户；人数已满则显示"已满员"
\end{itemize}

\subsubsection{业务规则}

\begin{itemize}
    \item 约练信息有效期 7 天，到期自动关闭
    \item 每人同时最多发布 3 条约练信息
    \item 确认后取消需提前 2 小时，频繁取消者限制发布权限
    \item 完成后系统自动触发双向评价邀请
    \item 搭子关系基于双方约练完成后互相添加，非单方面关注
\end{itemize}

\subsection{M5：成长档案模块}

\begin{table}[H]
\centering
\begin{tabularx}{\textwidth}{L{3cm}X C{0.8cm}}
\toprule
\textbf{功能} & \textbf{描述} & \textbf{级别} \\
\midrule
训练打卡 & 每次练舞后记录舞种、时长、地点 & P1 \\
学习数据统计 & 学舞天数、累计时长、课程数、舞种数 & P1 \\
成长时间线 & 以时间线展示学舞历程和关键节点 & P1 \\
阶段作品记录 & 上传练舞视频/照片作为阶段成果 & P2 \\
学习目标设置 & 设置周/月训练目标并追踪完成度 & P2 \\
成长报告 & 按月/季度生成学习报告 & P3 \\
成就徽章 & 达成里程碑解锁徽章（如连续打卡 30 天） & P3 \\
\bottomrule
\end{tabularx}
\end{table}

\subsubsection{成长数据指标}

\begin{table}[H]
\centering
\caption{成长数据展示设计}
\begin{tabularx}{\textwidth}{L{3.5cm}L{4cm}X}
\toprule
\textbf{指标} & \textbf{说明} & \textbf{展示形式} \\
\midrule
累计学舞天数 & 首次打卡至今 & 数字 + 日历热力图 \\
总训练时长 & 累计练舞小时数 & 数字 + 趋势折线图 \\
已上课程数 & 通过平台记录的课程 & 数字 \\
尝试舞种数 & 练习过的不同舞种 & 数字 + 舞种标签 \\
连续打卡天数 & 最长连续记录 & 数字 + 进度条 \\
本周/月训练量 & 当前周期训练时长和次数 & 柱状图 \\
\bottomrule
\end{tabularx}
\end{table}

\subsection{M6：社区与活动模块}

主要功能包括动态发布（图文/视频）、互动（点赞、评论、转发）、话题标签（\#韩舞推荐 \#零基础日记 \#约练日常），以及 Workshop 管理和活动日历。

\subsubsection{Workshop 业务流程}

\begin{enumerate}
    \item 舞室/主办方发布 Workshop 信息（老师、时间、地点、价格、人数上限、报名截止时间）
    \item 平台审核通过后上线，推送通知给关注相关舞种/老师的用户
    \item 用户浏览详情，查看老师介绍、往期评价、场地信息
    \item 用户选择场次，完成报名并支付
    \item 报名满员后，后续用户进入候补队列；有人取消时自动通知递补
    \item 活动当天用户到场扫码签到
    \item 活动结束后，系统邀请已签到用户进行评价
\end{enumerate}

\subsection{M7：商家管理模块}

主要功能包括舞室认领（提交营业执照等资质审核）、信息管理、课表管理、教练管理、预约管理、评价管理（申诉）、数据看板和 Workshop 管理。

\subsubsection{角色权限矩阵}

\begin{table}[H]
\centering
\caption{商家端角色权限}
\begin{tabularx}{\textwidth}{L{2.8cm}C{2cm}C{2cm}C{2cm}C{2cm}}
\toprule
\textbf{功能} & \textbf{舞室管理员} & \textbf{全职教练} & \textbf{签约教练} & \textbf{自由教练} \\
\midrule
编辑舞室信息 & $\checkmark$ & -- & -- & -- \\
管理课表 & $\checkmark$ & 仅自己 & 仅自己 & -- \\
发布 Workshop & $\checkmark$ & 需审批 & 需审批 & 独立发布 \\
查看舞室数据 & $\checkmark$ & 部分 & -- & -- \\
管理教练账号 & $\checkmark$ & -- & -- & -- \\
回复用户评价 & $\checkmark$ & 仅自己 & 仅自己 & 仅自己 \\
课程收益归属 & 舞室 & 舞室 & 按协议 & 教练本人 \\
\bottomrule
\end{tabularx}
\end{table}

\subsection{工程复杂性说明}

本项目表面上是信息平台，但工程复杂度并不低。它是包含定位搜索、结构化筛选、垂直评价体系、课程管理、约练匹配、成长记录和内容沉淀的复合型产品。不同模块的数据关系较强，既要支持用户从搜索到决策的快速路径，也要支持其长期使用中的学习和社交行为。真正的难点在于如何用一条用户路径把这些模块串联起来，而非简单堆叠功能。

% ============================================================
\section{用户价值与商业模式}
% ============================================================

\subsection{用户价值}

\begin{itemize}
    \item \textbf{对普通用户}：降低信息搜索和课程决策成本，不必在多个平台间跳转；通过搭子机制和成长记录降低中途放弃概率
    \item \textbf{对舞室和教练}：精准触达有学习意愿的用户，通过结构化展示和真实评价获得适合自身定位的学员，而非泛流量
\end{itemize}

\subsection{核心交易}

平台以更透明的决策信息、更稳定的练习陪伴和更清晰的成长反馈，换取用户的时间、信任和后续付费意愿。

\subsection{商业化路径}

\begin{table}[H]
\centering
\caption{商业化收入来源}
\begin{tabularx}{\textwidth}{L{2cm}L{3.5cm}X}
\toprule
\textbf{收入侧} & \textbf{收入形式} & \textbf{说明} \\
\midrule
舞室侧 & 推荐位/广告位 & 舞室付费获得搜索结果优先展示 \\
 & 试听投放 & 舞室在平台投放试听课吸引新学员 \\
 & 课程合作分成 & 通过平台完成的课程交易抽取佣金 \\
\midrule
用户侧 & 会员服务 & 解锁高级筛选、无限收藏、详细成长报告 \\
 & 成长工具增值 & 高级数据分析、成长报告导出 \\
 & 活动服务 & Workshop 报名手续费 \\
\midrule
联动合作 & 品牌合作 & 舞蹈用品、训练营、赛事合作 \\
\bottomrule
\end{tabularx}
\end{table}

平台连接的是线下真实消费场景而非低频浏览，通过在舞室选择、课程转化和用户留存三方面建立价值，具备从信息平台延伸到服务平台的现实可能。

\subsection{替换成本控制}

\begin{itemize}
    \item \textbf{用户侧}：无需放弃现有平台使用习惯，前期只在决策环节建立价值即可。用户即使不发内容、不约搭子，也可以先通过搜索和筛选获得收益
    \item \textbf{商家侧}：前期仅需更准确的信息展示和目标用户触达，后续逐步扩展到课表维护、试听投放和活动发布
    \item \textbf{推广策略}：聚焦单城市/区域，依赖内容传播（试听体验、避坑指南、约练记录）和供给端合作（舞室、校园舞社、活动主办方）而非大规模买量
\end{itemize}

% ============================================================
\section{MVP 验证方案}
% ============================================================

\subsection{MVP 范围}
BitDance 的 MVP 阶段将聚焦服务浅试型和行动型用户，帮助他们降低决策成本、提高留存。观望型用户通过内容社区逐步转化。首版围绕最核心的用户闭环收敛，优先实现：

\begin{itemize}
    \item 附近舞室搜索与舞种/课程筛选
    \item 结构化评价
    \item 试听收藏
    \item 约搭子
    \item 成长打卡
\end{itemize}

这样既能验证用户是否需要更垂直的决策平台，也能验证约练和记录功能是否足以延长使用周期。

\subsection{试点策略}

先在一个城市或高校周边商圈内试点，优先覆盖年轻用户密集、舞室供给较多的区域。目标对象以大学生和年轻上班族为主，因为他们既是舞蹈消费主力，也是最容易形成内容和社交关系的群体。

\subsection{成功指标}

优先验证三个核心假设：

\begin{enumerate}
    \item 用户是否愿意在平台上做舞室和课程决策
    \item 用户是否愿意留下评价和成长记录
    \item 用户是否愿意通过平台寻找搭子和约练
\end{enumerate}

三项跑通后，后续内容和商业化才有基础。

\subsection{为什么必须先做 MVP}

产品最大的风险不是能不能做出页面，而是会不会因功能过多而失焦。先用小范围 MVP 验证用户最痛的决策环节和最真实的陪伴需求，才能判断后续哪些模块值得扩展，哪些只适合作为增强功能保留。

% ============================================================
\section{风险与应对}
% ============================================================

\begin{table}[H]
\centering
\caption{风险分析与应对策略}
\begin{tabularx}{\textwidth}{C{0.6cm}L{3cm}X}
\toprule
\textbf{编号} & \textbf{风险} & \textbf{应对策略} \\
\midrule
R1 & 功能范围过大 & 先聚焦决策和约练主线，内容社区与商家后台分阶段推进，不在第一阶段做过重 \\
\midrule
R2 & 商家信息维护难 & 采用平台初始录入 + 商家认领完善 + 用户纠错补充三层机制 \\
\midrule
R3 & 评价体系被刷单 & 评价权重分层、结构化维度为主、异常识别风控；舞室仅可申诉不可删评 \\
\midrule
R4 & 舞室与教练合作关系复杂 & 通过角色划分（舞室/全职/签约/自由教练）、权限设计和交易留痕降低纠纷风险 \\
\midrule
R5 & 社交模块冷启动 & 约搭子建立在课程、舞种、地点等已有关系之上，让社交从学习场景自然生长 \\
\midrule
R6 & 平台沦为普通内容社区 & 始终坚持舞蹈学习主线，内容和社交功能均服务于核心路径，不做无关内容堆积 \\
\bottomrule
\end{tabularx}
\end{table}

% ============================================================
\section{立项建议}
% ============================================================

综合用户需求真实性、市场空白、替代方案局限和工程可行性，本项目具备明确的立项价值。

它解决的不是单纯找一家舞室的问题，而是舞蹈学习者从开始学舞到持续练舞过程中最关键的几步：\textbf{怎么选、怎么学、怎么坚持、怎么看到自己的进步}。只要这条主路径做扎实，平台就能从舞室查询工具逐步成长为面向舞者的学习与约练平台。

\vspace{0.5em}
\textbf{建议立项}，采用分阶段推进方式：
\begin{enumerate}
    \item \textbf{第一阶段}：舞室决策和约练闭环（MVP）
    \item \textbf{第二阶段}：内容社区与活动合作
    \item \textbf{第三阶段}：商家服务与商业化能力
\end{enumerate}

% ============================================================
\section*{参考资料}
\addcontentsline{toc}{section}{参考资料}
% ============================================================

\begin{enumerate}[label={[\arabic*]}]
    \item 国家统计局：2025 年居民收入和消费支出情况。\\\url{https://www.stats.gov.cn/sj/zxfb/202601/t20260119_1962321.html}
    \item 国家统计局：2025 年全民健身活动状况调查统计调查制度。\\\url{https://www.stats.gov.cn/fw/bmdcxmsp/bmzd/202506/t20250605_1960060.html}
    \item 中国舞蹈家协会街舞委员会：《中国街舞产业发展研究报告》。\\\url{https://www.chhuc.org/18158.html}
    \item 国家体育总局办公厅：关于政协第十四届全国委员会第一次会议第 04245 号提案答复的函。\\\url{https://www.sport.gov.cn/bgt/n25546387/c27068331/content.html}
    \item 第一财经：韩流奔涌全球的文化视角解因。\\\url{https://m.yicai.com/news/102227235.html}
    \item CNNIC：第 56 次中国互联网络发展状况统计报告。\\\url{https://www.cnnic.net.cn/NMediaFile/2025/0730/MAIN1753846666507QEK67ZS9DH.pdf}
\end{enumerate}

\end{document}
